\paragraph{Source stress and the test-clock domain.}
The autonomous wave construction determines its field force through
$hG_a+G_ah=2(D_aK_q+K_qD_a)$, with $h^2=K_q$ and
$D_a=\operatorname{diag}(F_a)$. The complete finite Ehrenfest equation
retains the oscillator boundary term; geometry, native and field energy
rates sum to zero. Full stress and momentum instruments retain all
complementary outcomes. Interval certificates for all64 native vacuum
sectors show a frequency-ratio spread of
$[87.8991138,87.8991463]$ppm among the seven dressed native clocks.
A gap/resolvent theorem gives
$|(f(s)-f(0))/s-\langle e^{Q_1}\rangle_0|
\leq[(17/250)^2/(9/50)]s$ on the stated probe domain. Thus the common
test-clock response has a quantitative backreaction bound, while exact
finite-strength universality fails in this construction. The source
choices and separate clock preparations are retained; physical $G$,
laboratory clocks and an Einstein stress tensor are not inferred.
See \texttt{model/openwave\_dynamic\_stress}.
